\section{Restricciones, herramientas y Q\&A de clase: conectar prototipos de investigación con sistemas reales}

La clase no terminó cerrando con cifras SOTA, sino que se centró en discutir qué interfaces aún no están resueltas: horizonte corto, único nivel temporal, entornos juguete (toy environments), imagen objetivo (image goal), robots reales, arquitectura de agente y destilación de políticas. El Q\&A también reiteró constantemente que JEPA es un marco de entrenamiento, no un nuevo tipo de red neuronal, y que el modelo del mundo (world model) puede combinarse con difusión, Transformer, política RL y control clásico.

\subsection{Cuatro limitaciones públicas corresponden a cuatro líneas de investigación}

La diapositiva 53 convierte la limitación en oportunidad de investigación. El horizonte corto exige un rollout más estable o corrección en bucle cerrado; el único nivel temporal requiere estado jerárquico; el montaje juguete (toy setup) exige percepción real y cambio de dominio (domain shift); una imagen objetivo irrealista requiere codificador de lenguaje/tarea o interfaz de recompensa. Ninguna de ellas se resuelve simplemente ``ampliando el modelo''.

\lecturefigure{slide-053.jpg}{Limitaciones de LeWorldModel: horizonte corto, escala única, entorno juguete e interfaz de objetivo}{CS25 V6 oficial, p.~53}

\begin{knowledgebox}{De la limitación al módulo de sistema}
\begin{tabularx}{\textwidth}{L{0.24\textwidth} X}
\toprule
Limitación & Capacidad del sistema a añadir \\
\midrule
planeación a corto plazo & rollout consciente de incertidumbre, replaneación, híbrido política/modelo-mundo \\
nivel temporal único & latente jerárquico, abstracción de eventos, memoria a largo plazo \\
entornos juguete & percepción robusta, simulación-a-realidad, observabilidad parcial, restricciones de seguridad \\
imagen objetivo irrealista & codificador de lenguaje/tarea, modelo de recompensa, especificación de restricciones \\
\bottomrule
\end{tabularx}
\end{knowledgebox}

\teachervoice{00:55:14--00:57:15, Lucas identificó explícitamente estos proyectos como limitaciones u oportunidades de investigación: los planificadores actuales son principalmente a corto plazo y con representación temporal de un solo nivel, y las interfaces de entorno y objetivo siguen estando lejos de la robótica abierta.}

\subsection{stable-worldmodel: la comparabilidad reproducible también es contribución metodológica}

La diapositiva 54 muestra una librería unificada que coloca LeWM, PLDM, DINO-WM y otras líneas base (baselines), junto con solvers como CEM, MPPI y Adam, y entornos como Push-T, TwoRoom, OG-Bench y DMC en una misma interfaz. Su valor científico radica en reducir la incomparabilidad que causan el preprocesamiento de datos y las implementaciones de planificadores distintas en cada artículo.

\lecturefigure{slide-054.jpg}{stable-worldmodel: capa experimental unificada de baselines, solvers y entornos}{CS25 V6 oficial, p.~54}

\teachervoice{00:57:15--00:59:16, el docente invitó a los estudiantes a usar y contribuir con ``stable-worldmodel'', enfatizando que las baselines, solvers, entornos, pruebas y documentación ya están unificados. El punto clave aquí es la reproducibilidad, no convertir una librería en el único estándar.}

\begin{importantbox}{Una API unificada no sustituye a un protocolo unificado}
La misma librería aún requiere fijar resolución de observación, normalización de acciones, división del conjunto de datos (dataset split), presupuesto del planificador, semillas aleatorias y definición de éxito. La API resuelve la interfaz de código, pero el protocolo de benchmark es lo que garantiza la comparabilidad científica; ambos son indispensables.
\end{importantbox}

\subsection{Physical AI: por qué la predicción de consecuencias de acción es un objetivo de entrenamiento independiente}

En el Q\&A, Lucas expresó escepticismo sobre si los VLA actuales poseen una comprensión del mundo confiable: si el modelo no se entrenó para predecir las consecuencias de sus propias acciones, puede funcionar bien en tareas de imitación familiares, pero carecería de verificación explícita para nuevas interacciones. Hazel añadió que el modelo del mundo también puede servir como evaluador de políticas. Esta perspectiva debe considerarse como un juicio del orador, no como una conclusión empírica universal para todos los VLA.

\teachervoice{00:59:16--01:01:17, Lucas opinó que si el agente físico no puede predecir las consecuencias de la acción, difícilmente podrá interactuar de manera confiable; Hazel complementó que el modelo del mundo puede evaluar una política candidata. Las notas recuerdan que esto es un juicio prospectivo que requiere validación con experimentos en robots reales.}

\begin{warningbox}{VLA y modelo del mundo no son categorías mutuamente excluyentes}
Un VLA puede contener internamente un objetivo predictivo, un simulador latente o un valor/modelo del mundo; el modelo del mundo también puede implementarse con Transformers de lenguaje y visión. La verdadera diferencia radica en la interfaz de entrenamiento: si el modelo aprende explícitamente $p(s_{t+1}\mid s_t,a_t)$ y si, durante el despliegue, usa esa predicción para verificar o buscar acciones.
\end{warningbox}

\subsection{¿Es necesaria la máscara? Lo necesario es eliminar atajos (shortcuts)}

Un estudiante preguntó si la pérdida de predicción era insuficiente. La respuesta de Hazel fue prudente: para cualquier modelo del mundo, la máscara no es el único método lógicamente posible; pero en un montaje centrado en objetos (object-centric setup), predecir sin oclusiones facilita aprender las dinámicas propias del objeto y no garantiza las dinámicas basadas en interacción. La máscara es un medio para sacar los atajos fuera del espacio de soluciones factibles, no una ley universal.

\teachervoice{01:01:17--01:03:15, respondiendo a ``¿máscara necesaria?``, el docente aclaró que la pérdida de predicción por sí misma puede aprender ciertos modelos del mundo, pero podría limitarse a aprender solo el movimiento del objeto; la máscara de objetos refuerza las dinámicas de interacción y no debe escribirse como la única receta obligatoria para todos los sistemas.}

\begin{importantbox}{Pregunta general sobre el diseño de objetivos auto-supervisados}
No se deba preguntar solo ``¿puede bajar la pérdida?``, sino: \textbf{cuál es el atajo más barato?} Si la tarea objetivo requiere relaciones, memoria a largo plazo o consecuencias de acción, el entrenamiento debe eliminar sistemáticamente las vías que pueden completarse solo con estadísticas locales, fuga de identidad o correlaciones estáticas.
\end{importantbox}

\subsection{Diferencias de interfaz entre agente JEPA-native y agente LLM con herramientas}

Lucas describió al agente JEPA-native como un modelo predictivo condicionado a la acción más control clásico: ingresa el estado actual y la acción, predice el futuro y luego realiza búsqueda sin ejemplos usando MPC. Los agentes LLM aprenden habitualmente a usar herramientas mediante bucles de texto/citación de herramienta, y no necesariamente aprenden explícitamente las transiciones de estado de acciones físicas continuas. Ambos pueden combinarse; por ejemplo, un modelo de lenguaje descompone el objetivo y el modelo del mundo se encarga de las consecuencias de acción de nivel inferior.

\teachervoice{01:03:19--01:05:20, la clase enfatizó que la diferencia del agente JEPA-native radica en ser un modelo predictivo condicionado a la acción, no en usar algún backbone nuevo. Una vez aprendido el comportamiento dinámico (dynamics), se puede conectar directamente con MPC sin necesidad de post-entrenar cada objetivo como una política primero.}

\begin{knowledgebox}{Capas combinables del stack de agente}
El modelo de lenguaje puede encargarse de la explicación de tareas y subobjetivos de alto nivel; el modelo del mundo predice el estado continuo del entorno; el planificador/solver busca acciones; la política ejecuta rápido; el verificador y escudo de seguridad gestionan las restricciones. Separando estas capas, se puede localizar si el fallo proviene de semántica, dinámica, optimización o ejecución.
\end{knowledgebox}

\subsection{JEPA no es el antónimo del Transformer}

La pregunta ``¿JEPA genera menos alucinaciones que el Transformer?'' mezcla marco y arquitectura. Lucas aclaró explícitamente que los codificadores y predictores de LeWorldModel son en sí mismos Transformers; JEPA describe qué se predice durante el entrenamiento, mientras que Transformer describe cómo interactúan los tokens. El modelo del mundo seguirá generando futuros erróneos debido a falta de datos, insuficiencia de capacidad, error del modelo o acciones fuera de límites detectadas por el controlador.

\teachervoice{01:05:20--01:07:22, Lucas corrigió la confusión de categorías: JEPA no es una nueva arquitectura, sino un marco para aprender modelos del mundo. Incluso usando JEPA, datos malos o modelos malos seguirán alucinando; simplemente los tipos de error no son exactamente los mismos que la generación de texto sin grounding de los LLM.}

\begin{warningbox}{Cuatro fuentes de alucinación en el modelo del mundo}
\begin{enumerate}
  \item El codificador pierde estados relevantes para la decisión;
  \item El predictor acumula errores en horizontes largos;
  \item La distribución de entrenamiento carece de interacciones clave;
  \item El planificador propone acciones inválidas o fuera de límites del entorno.
\end{enumerate}
Por ello, ``más grounded'' es una hipótesis comprobable, no una propiedad garantizada automáticamente por el nombre JEPA.
\end{warningbox}

\subsection{Difusión, MPC y destilación de políticas pueden coexistir}

Otra pregunta trataba a los modelos de difusión y al modelo del mundo como productos competidores. El ponente señaló que la difusión puede servir como predictor del modelo del mundo, propuesta de acción o política; el modelo del mundo es un rol funcional, no una familia específica de modelos probabilísticos. Para horizontes largos, el sistema también puede usar MPC costoso para generar comportamientos de alta calidad y luego destilarlos en una política rápida: las situaciones familiares siguen rutas reflejas y las difíciles retroceden a la planificación.

\teachervoice{01:07:22--01:10:59, la clase dio dos respuestas de sistema: la difusión puede integrarse en el pipeline JEPA/modelo del mundo; los futuros agentes podrían funcionar como sistemas rápidos/lentos, usando directamente la política en casos simples y llamando a MPC en casos difíciles, mientras se destila la experiencia de planificación de vuelta a la política.}

\begin{importantbox}{Ciclo cerrado de ingeniería entre política rápida y planificador lento}
El modelo del mundo offline aprende las dinámicas; el MPC online genera acciones para nuevos objetivos; las trayectorias exitosas entran en el buffer de reproducción (replay buffer); la destilación de políticas aprende un mapeo rápido; estados de baja confianza o fuera de distribución reactivan el planificador. Este ciclo conecta modelo, control y recolección continua de datos, pero introduce problemas de sesgo del modelo, deriva de la política y verificación de seguridad.
\end{importantbox}

\subsection{Resumen de la sección}

Los sistemas reales no elegirán entre JEPA, Transformer, difusión, MPC o políticas en cinco categorías. Cada uno responde a objetivos de entrenamiento, estructura de red, modelado probabilístico, búsqueda online y ejecución rápida. Las principales lagunas en los prototipos de investigación actuales son horizonte largo, jerarquía, entornos reales e interfaces de objetivo; el benchmarking unificado mejora la calidad de la comparación, pero no sustituye al despliegue real ni a las pruebas de seguridad.
